% TOP_PROBLEM: {"id": 3082, "title": "Inequality of the means", "category_id": 6, "category": {"id": 6, "name": "geometry", "display_name": "Geometry", "description": "Euclidean and non-Euclidean geometry, geometric structures.", "slug": "geometry", "order_index": 6, "created_at": "2026-07-31T15:26:25.671Z"}, "status": "open", "problem_number": "OPG-36901", "set_id": 12, "statusQualification": "Ordinary congruent packing allows rotations; the stronger axis-parallel variant is identified rather than silently imposed."}
% FIRST_POSTED: 2026-10-01
% ATTEMPT_STATUS: unresolved
% COMPLETION: 5
% MODEL: GPT 6 Astra Ultra
% UnsolvedMath ID 3082; source catalogue code OPG-36901
% !TeX program = lualatex
% Research attempt, 2026-10-01. Canonical statement preserved verbatim.
\documentclass[11pt,a4paper]{article}
\usepackage[margin=25mm]{geometry}
\usepackage{fontspec}
\setmainfont{lmroman10-regular.otf}[BoldFont=lmroman10-bold.otf,
 ItalicFont=lmroman10-italic.otf,BoldItalicFont=lmroman10-bolditalic.otf]
\usepackage{amsmath,amssymb,mathtools}
\usepackage{unicode-math}
\setmathfont{latinmodern-math.otf}
\AtBeginDocument{\let\setminus\smallsetminus}
\usepackage{xurl,hyperref}
\hypersetup{colorlinks=true,urlcolor=blue}
\setcounter{secnumdepth}{0}
\setlength{\parindent}{0pt}
\setlength{\parskip}{5pt}
\setlength{\emergencystretch}{3em}
\begin{document}
\section*{Inequality of the means}\label{3082}
{\small UnsolvedMath ID 3082 · OPG-36901 · Geometry · OpenGarden}
\subsection{Definitions and mathematical statement}
Let $n\ge1$ and let $a_1,\ldots,a_n$ be arbitrary positive real numbers.
Set $s=\sum_{i=1}^n a_i$, let
$B=\prod_{i=1}^n[0,a_i]$ be a rectangular box, and let $Q=[0,s]^n$.
A packing of $m$ congruent copies of $B$ in $Q$ is a family of Euclidean
isometries $T_1,\ldots,T_m$ such that $T_j(B)\subseteq Q$ for every $j$
and the interiors of different copies are disjoint. Boundary contacts
are allowed; boxes may be moved and rotated but not cut or deformed.

Does such a packing always exist with $m=n^n$? The quantifiers range
over every dimension and every choice of side lengths. Permuting the
side lengths changes only the presentation of the same box, not the
available shape. The definition uses the ordinary unrestricted meaning
of geometric packing in the catalogue question. Requiring every edge to
be parallel to a coordinate axis would be a stronger version; the source's
illustrated constructions use such orientations.

A necessary volume comparison is
\[
 n^n\prod_{i=1}^n a_i\le\left(\sum_{i=1}^n a_i\right)^n,
\]
which is exactly the arithmetic--geometric mean inequality after taking
$n$th roots. The problem asks for actual disjoint placement, not just this
numerical inequality. When all lengths equal $a$, the $n\times\cdots\times n$
grid tiles the target cube. In the unequal case, unused space is permitted
and the copies need not cover the entire cube.
\subsection{Short English statement}
Can $n^n$ congruent boxes with side lengths $a_1,\ldots,a_n$ always fit into a cube of side $a_1+\cdots+a_n$?
\subsection{Research attempt}
\paragraph{Scope and literature check.}
The constructions on Bar-Natan's page [S3] already establish the two- and
three-dimensional cases and explain composition of dimensions. This attempt
gives an explicit coordinate certificate for the two-dimensional construction,
then extends it to an even-dimensional family of side lengths. These are
restricted results; the necessary AM--GM volume inequality is never used as
if it supplied a packing. Sources [S2,S3] were checked on 1 October 2026.

\paragraph{A packing with a visible remainder.}
Assume $a\ge b>0$ and put $s=a+b$. In $[0,s]^2$ use the four rectangles
\[
\begin{split}
 R_1&=[0,a]\times[0,b],& R_2&=[a,s]\times[0,a],\\
 R_3&=[b,s]\times[a,s],& R_4&=[0,b]\times[b,s].
\end{split}
\]
Each has side lengths $a,b$, in one of the two permitted orientations.
Each pair has disjoint interiors: consecutive rectangles are separated by
$x=a$, $y=a$, $x=b$, or $y=b$ as appropriate; the two remaining pairs are
separated by the horizontal or vertical gap from $b$ to $a$. All endpoints
are in $[0,s]$. The uncovered interior is precisely $(b,a)^2$, giving
$s^2-4ab=(a-b)^2$. At $a=b$ the gap disappears and the construction becomes a
tiling. Swapping the names of the sides handles $b>a$.

\paragraph{An even-dimensional restricted theorem.}
Let $n=2k$. Suppose the side lengths can be paired
$(a_{2j-1},a_{2j})$, $1\le j\le k$, so that each pair has the same sum $t$.
Then $s=kt$. Form the Cartesian product of the four-rectangle packing for
each pair. It gives $4^k$ disjoint $2k$-boxes inside $[0,t]^{2k}$, each
congruent to the requested box; exchanging coordinates independently within
a pair is an orthogonal transformation in the full space.
Now partition $[0,kt]^{2k}$ into $k^{2k}$ translates of $[0,t]^{2k}$ and
place that product packing in every small cube. Interiors from distinct
cubes cannot intersect, so the total is exactly
\[
 4^k k^{2k}=(2k)^{2k}=n^n.
\]
This works for arbitrarily unequal lengths within each pair and arbitrarily
large even dimension, provided the pair sums agree. Its occupied volume is
$n^n\prod_i a_i$ as required; the geometric construction, rather than the
volume calculation, proves feasibility.

\paragraph{Why the most direct generalization fails.}
For unequal pair sums $t_j$, the same product has container side lengths
$t_1,t_1,\ldots,t_k,t_k$. Repeating it on a grid in a cube of side
$s=\sum_jt_j$ guarantees only
\[
 M=4^k\prod_{j=1}^k\left\lfloor s/t_j\right\rfloor^2
\]
copies. There is no universal inequality $M\ge (2k)^{2k}$. For $k=2$,
$t_1=3,t_2=2$, one obtains $M=16\cdot1^2\cdot2^2=64$, whereas $4^4=256$.
For example the four side lengths $(1,2,1,1)$ realize these pair sums.
This does not contradict the known four-dimensional theorem: it disproves
only this particular fixed-product grid argument. Mixed orientations of
larger blocks, absent from that argument, can change the answer.

\paragraph{Verification and remaining gap.}
The accompanying finite checks verify containment and all pairwise interior
separations for integer pairs $1\le b\le a\le8$, including the equal-side
boundary case, and verify the count identities. The proof itself covers
all positive real lengths. To address general $n$ one still needs a placement
that accommodates unequal block sums without losing the required number of
copies. No exhaustive search in dimension five is claimed.

\paragraph{Reproducibility.}
The finite calculations above are implemented in
\texttt{scripts/check\_attempt\_batch\_2026\_10\_01.py}; run it with Python 3
from the repository. It uses exact integer arithmetic and no external packages.

\paragraph{Final status.}
\textbf{UNRESOLVED.} There is an explicit packing theorem for equal-sum pairs,
and a numerical obstruction to the naive unequal-sum extension. These are
not a resolution for arbitrary side lengths and dimensions. Completion
estimate: 5\%.

\subsection{Sources}
\begin{itemize}
\item[S1] ULAM.AI and the UnsolvedMath Contributors, supplied problems.json, record 3082 (OPG-36901), OpenGarden collection. Catalogue identity and source transcription; CC BY 4.0.
\url{https://huggingface.co/datasets/ulamai/UnsolvedMath}.
\item[S2] Open Problem Garden, Inequality of the means. Original problem and discussion; the mathematical formulation below is expanded from the supplied source record.
\url{http://www.openproblemgarden.org/op/inequality_of_the_means}.
\item[S3] D. Bar-Natan, Inequality of the Means, author’s illustrated formulation and constructions.
\url{https://www.math.toronto.edu/~drorbn/projects/ArithGeom/}.

\end{itemize}
\end{document}
